\head{第十七版来源核验与合并结论}
本版恢复并检查第十六版核心续作包与原工程包，两者的现行源码和构建说明共有132份可比文件逐字相同。实际第十六版进度为第0章止于（23.2.10，ii），第V卷第一篇止于§6。任何与这一源码不一致的旧对话进度均未作为完成证据。

\audit{D17-I01}{Dot实存译稿及原件追溯}{56批已纳入来源清单；不把归档成功等同于内容审定}
Dot续作总包有54批EGA V独立工程，本版逐批取得实际TeX、进度及相应原件；总包之后另恢复两批第16节工程。54批按内容顺序分为§§7--10的18批、§§12--13的19批、§§14--15的17批。第16节两批与主稿短篇重合，作为同文校订处理。

8.14--8.15批次中的\texttt{前批8.13.tex}为备份，不作为第二份新增正文。Dot总包内另有其他研究项目资料，原包整体归档以保留交付，但这些无关资料没有进入EGA的正文或完成度统计。每批原始文件名、哈希、原进度字段及处理位置见工程的来源映射；过时进度作为历史记录保存，不覆盖本版接续点。

\audit{D17-I02}{实际漏段与本轮新译}{两处据底本补齐，不计作Dot既有成果}
第11节原文位于\texttt{EGAVaiello.pdf}第41--43页，位于Dot已译（10.4）与§12之间；本版补译全节及原脚注57--60。第14节第54--55页相对锥投影至扩展锥投影的段落在原包中漏接，本版补译并保留原脚注67和交换图。两处均按实际底本逐段处理，具体说明见相应核查条目。

\audit{D17-I03}{正文纠正与原文留证}{确定修订进入正文，无法修复的反例项明确标示，未决问题仍待核}
本版将原分批稿中已指出却尚留在正文的若干错误迁入正确表述。处理包括：纠正（8.4.1）假设及其局部维数证明、（10.1）可构造例外集与闭包的区别、§11的Grassmann构造与有限性条件、§12退化层维数的适用条件、第14节跨页图式的基底以及旧短篇的类型与Jacobian判据问题。

对于（7.4）和命题14.4(a)的一条结论，现有反例足以否定所印假设，但尚未恢复完整有效的修订版；正文因此说明不能采用，原陈述或公式及理由保存在脚注和报告。没有以猜测补齐原缺文，也没有把英译转录错误直接归因于尚未核对的法文手稿。

\audit{D17-I04}{构建、保留与接续}{以本次实际执行记录为准}
本版使用\texttt{ctexbook}默认版式，主稿和累计报告由源码重新构建；单文件TeX由同一份分文件工程展开。日志检查、独立编译、页面渲染与源码保留比对的最终结果载于工程\texttt{checks}目录及\texttt{PROGRESS.md}。旧日志和Dot分批自述属于历史资料，不能替代这些本次记录。

新工程继续保留全部旧译文；其中旧短篇一份有明确登记的同文校订和修正，原版本另存档。旧核查正文保持，文件保留核验不等于重新证明其中的数学结论。后续真实接续点为第IV章内容提要、章首与（1.1.1），以及第V卷英文稿第66页的交换图。

